\documentclass{article}

\usepackage{amsmath, amssymb}
\usepackage{graphicx}
\usepackage[
    left=0.6in,
    right=0.6in,
    top=0.6in,
    bottom=0.6in
]{geometry}

\setlength{\parindent}{0pt}
\setlength{\parskip}{3pt}
\linespread{0.96}

\title{RM -- Credit Instruments and Risk Homework}
\author{Yueqi Lin, Yajie Zhang, Mat Zaga}
\date{May 2026}

\begin{document}

\maketitle

\section{Exercise 1}

\textbf{Consider a corporate bond:}
\begin{itemize}
    \item with notional $N$ (compare to \$1 in the lecture notes);
    \item issued by a firm with constant default intensity $\lambda$ (not a function of time);
    \item final maturity $T$, at which the notional is returned to the bond holders;
    \item constant continuous coupon $c$, i.e. in any time interval $[t, t+dt]$, where $t < T$, the bond holders are paid $N \cdot c \cdot dt$ as long as there is no default prior to $t$;
    \item recovery rate $R$, i.e. at the time of default (if before $T$), the bond holders receive $N \cdot R$.
\end{itemize}

Assume the risk-free rate is constant at $r > 0$.

\begin{align*}
\mathbb{P}(\tau > t) &= e^{-\lambda t}, \\
\mathbb{P}(\tau \in dt) &= \lambda e^{-\lambda t} dt
\end{align*}

%--------------------------------------------------

\section*{Case 1: $R = c = 0$}

\textbf{Question:} What is the bond value at $t = 0$?

\vspace{4pt}

\textbf{Solution:}

\begin{align*}
V_0 &= N e^{-rT} \mathbb{P}(\tau > T) \\
    &= N e^{-rT} e^{-\lambda T} \\
    &= N e^{-(r+\lambda)T}
\end{align*}

\[
\boxed{V_0 = N e^{-(r+\lambda)T}}
\]

%--------------------------------------------------

\section*{Case 2: $R = 0,\ c \neq 0$}

\textbf{Question:}
\begin{itemize}
    \item[(a)] What is the value of the coupon stream at $t = 0$?
    \item[(b)] What is the total value of the bond at $t = 0$?
    \item[(c)] What is the limit of the initial bond value as $T \to \infty$?
    \item[(d)] At which value of $c$ is the initial bond value equal to $N$ (par coupon)?
\end{itemize}

\vspace{4pt}

\textbf{Solution:}

\textbf{(a) Coupon value}

\begin{align*}
V_{\text{coupon}}
&= \int_0^T Nc \, e^{-rt} \mathbb{P}(\tau > t)\, dt \\
&= Nc \int_0^T e^{-rt} e^{-\lambda t} dt \\
&= Nc \int_0^T e^{-(r+\lambda)t} dt \\
&= \frac{Nc}{r+\lambda}\left(1 - e^{-(r+\lambda)T}\right)
\end{align*}

\textbf{(b) Total value}

\begin{align*}
V_0
&= V_{\text{coupon}} + V_{\text{notional}} \\
&= \frac{Nc}{r+\lambda}\left(1 - e^{-(r+\lambda)T}\right)
+ N e^{-(r+\lambda)T}
\end{align*}

\[
\boxed{
V_0 =
\frac{Nc}{r+\lambda}\left(1 - e^{-(r+\lambda)T}\right)
+ N e^{-(r+\lambda)T}
}
\]

\textbf{(c) Limit as $T \to \infty$}

\[
\boxed{
\lim_{T \to \infty} V_0 = \frac{Nc}{r+\lambda}
}
\]

\textbf{(d) Par coupon}

\begin{align*}
1
&= \frac{c}{r+\lambda}\left(1 - e^{-(r+\lambda)T}\right)
+ e^{-(r+\lambda)T}
\end{align*}

\[
\boxed{c = r + \lambda}
\]

%--------------------------------------------------

\section*{Case 3: $R \neq 0,\ c \neq 0$}

\textbf{Question:}
\begin{itemize}
    \item[(a)] What is the value of the recovery component at $t = 0$?
    \item[(b)] What is the total value of the bond at $t = 0$?
    \item[(c)] What is the limit of the initial bond value as $T \to \infty$?
    \item[(d)] At which value of $c$ is the initial bond value equal to $N$ (par coupon)?
\end{itemize}

\vspace{4pt}

\textbf{Solution:}

\textbf{(a) Recovery value}

\begin{align*}
V_{\text{recovery}}
&= \int_0^T NR \, e^{-rt} \lambda e^{-\lambda t} dt \\
&= NR \lambda \int_0^T e^{-(r+\lambda)t} dt \\
&= \frac{NR\lambda}{r+\lambda}\left(1 - e^{-(r+\lambda)T}\right)
\end{align*}

\textbf{(b) Total value}

\begin{align*}
V_0
&= V_{\text{coupon}} + V_{\text{notional}} + V_{\text{recovery}} \\
&= \frac{N(c + R\lambda)}{r+\lambda}\left(1 - e^{-(r+\lambda)T}\right)
+ N e^{-(r+\lambda)T}
\end{align*}

\[
\boxed{
V_0 =
\frac{N(c + R\lambda)}{r+\lambda}\left(1 - e^{-(r+\lambda)T}\right)
+ N e^{-(r+\lambda)T}
}
\]

\textbf{(c) Limit as $T \to \infty$}

\[
\boxed{
\lim_{T \to \infty} V_0 = \frac{N(c + R\lambda)}{r+\lambda}
}
\]

\textbf{(d) Par coupon}

\begin{align*}
1
&= \frac{c + R\lambda}{r+\lambda}\left(1 - e^{-(r+\lambda)T}\right)
+ e^{-(r+\lambda)T}
\end{align*}

\[
\boxed{c = r + \lambda(1 - R)}
\]

\end{document}